\begin{abstract}
When an LLM judge supplies the reward for a dialogue policy, it is usually shown the candidate
turn and nothing after it. In multi-turn goal-oriented dialogue that proxy is short-sighted: the
turn that reads best on its own is not necessarily the turn that leads anywhere. Motivational
interviewing (MI) makes the gap concrete, because the value of a counsellor's utterance lies in
what the patient says next. We present \textbf{\LAGRPO}: before the judge scores a candidate
therapist turn, the patient simulator and the current policy continue the conversation for $K$
further turns and the judge scores the continuation. The look-ahead reward comes from
preference-tree optimisation for MI; moved to Group Relative Policy Optimization (GRPO), it
reweights every sampled candidate instead of selecting a preference pair. In a controlled pair of
ten-iteration runs that differ only in the horizon, $K{=}5$ leads on the rewarded rubric from the
fourth iteration on, beats the $K{=}0$ policy's best checkpoint on that rubric as well as its
last, leads on all eight MI evaluation instruments under both the training oracle and a held-out
judge from another model family, and replicates on a fresh evaluation draw. The horizon also
decides what is learned. Turn-level reward induces over-praise, an MI-inconsistent behaviour whose
rate quadruples over the $K{=}0$ run and which a judge-free lexical marker corroborates; the
look-ahead policy asks questions instead, elicits more change talk, and under the held-out judge
drifts toward MI-inconsistency at under half the turn-level rate. Each arm is a single training
run, and at the winning checkpoint the training oracle's scores reach the ceiling of its scale, so
the held-out judge carries the per-conversation claims.
\end{abstract}
